\section{Derivation of Relevant Dynamics}

\subsection{Notation}

\begin{itemize}
  \item $\mathcal{I}$ denotes an intertial frame, with its $z$ axis pointing in the direction against gravity.
  \item The position vector in frame $\mathcal{A}$ to a point $Q$ relative to the frame origin is denoted as ${}^\mathcal{A}\mathbf{p}_{Q}$.
  \item $\rotmat{A}{B}$ denotes the rotation matrix from frame $\mathcal{B}$ to frame $\mathcal{A}$. This is such that ${}^\mathcal{A}\mathbf{p}=\rotmat{A}{B}{}^\mathcal{B}\mathbf{p}$.
  \item Let $[\mathbf{x}]_\times \in \mathbb{R}^{3\times 3}$ be the skew-symmetric matrix such that $[\mathbf{x}]_\times\mathbf{y} = \mathbf{x} \times \mathbf{y}$, where $\times$ is the cross product operator in $\mathbb{R}^3$.
\end{itemize}


\subsection{System Dynamics Formulation}

With the above notation in mind, we formulate the state of the vehicle as a whole as the following vector:

$$
\mathbf{x} = 
\begin{bmatrix}
  \mathbf{p} \in \mathbb{R}^3 \\
  R \in SO(3) \\
  \mathbf{v} \in \mathbb{R}^3 \\
  {}^\mathcal{B}\boldsymbol{\omega} \in \mathbb{R}^3 \\
  \boldsymbol{\omega}_w \in \mathbb{R}^4
\end{bmatrix}
$$
And the corresponding control input is also defined as:

$$
\mathbf{u} = \begin{bmatrix}
  \delta \in \mathbb{R} \\
  \boldsymbol{\tau} \in \mathbb{R}^2
\end{bmatrix}
$$

From this, we now derive the full state space model of the system. First, we define the translational and rotional kinematics of the car as a whole via Newton's second law:

\begin{equation}
  \begin{gathered}
    m\dot{\mathbf{v}} = m\mathbf{g} + \sum_i \mathbf{f}_i \\
    \dot{\mathbf{p}} = \mathbf{v}
  \end{gathered}
\end{equation}

Where $\mathbf{f}_i\in\mathbb{R}^3$ comes from the wheel dynmaics, detailed later. And similarly, for rotational kinematics:

\begin{equation}
  \begin{gathered}
    \mathbb{I}_G{}^\mathcal{B}\dot{\boldsymbol{\omega}} + {}^\mathcal{B}\boldsymbol{\omega} \times \left(\mathbb{I}_G{}^\mathcal{B}\boldsymbol{\omega}\right) = \boldsymbol{\tau}(\mathbf{x,u}) \\
    \dot{R} = R[{}^\mathcal{B}\boldsymbol{\omega}]_\times
  \end{gathered}
\end{equation}

And the formulation for the update of the orientation uses the built-in \mintinline{python}{jaxlie.SO3.exp} formulation, which evolves on $SO(3)$ while using a quaternion ($q\in S^3$) representation internally within the \mintinline{python}{jaxlie} library.

The quantity $\mathbf{f}_i$ is the force generated by the $i$-th wheel, and is calculated by both contact dynamics in $z$ and slip dynamics in both $x$ and $y$. To calulate the wheel dynamics, we first start wth the normal force $f_{z,i}$:

\begin{equation}
  f_{z,i} = k_n \cdot \mathrm{ReLU}(z_0-z) + c_n \cdot \mathrm{ReLU}(0-v_z)
\end{equation}

where $f_{z,i}$ is enforced to be positive, such that the car is not pulled into the ground (i.e. $z_0-z>0$). $z_0$ is the initial height, and $z$ is the current height of the wheel. $v_z$ is the vertical velocity of the wheel, and $k_n$ and $c_n$ are the stiffness and damping coefficients for the normal force, respectively.

%TODO: add reference for Pacejka slip model, and also add more details on the slip model itself, such as the stiffness coefficients and the maximum slip ratio.
Then, the wheel slip \textit{ratios} are then found via the Pacejka slip model, which is a commonly used empirical model for tire slip dynamics. The slip ratios are defined as follows:

\begin{equation}
  \begin{gathered}
    \kappa_i = \kappa_\mathrm{max} \cdot \frac{r\omega_i - v_{t,i}}{\sqrt{v_{t,i}^2+\epsilon^2}} \\
    \alpha_i = \mathrm{arctan}\left(\frac{v_{n,i}}{\sqrt{v_{t,i}^2+\epsilon^2}}\right) \\
  \end{gathered}
\end{equation}

In equation $(4)$, the tolerance $\epsilon$ is added to prevent division by zero when the tangential velocity $v_{t,i}$ is very small. Furthermore, $\kappa$ and $\alpha$ are also both bounded by a $\tanh$ function each, in order to bound large slip changes that may induce negative RPM values in the model. The slip ratios $\kappa_i$ and $\alpha_i$ are then used to calculate the longitudinal and lateral forces generated by the wheel, via $f_{x,i}^\star=C_\kappa \cdot \kappa$ and $f_{y,i}^\star=-C_\alpha \cdot \alpha$, where $C_\kappa$ and $C_\alpha$ are the longitudinal and lateral slip stiffness coefficients, respectively. From this, all forces are scaled by the maximum normal force, formulated as:

\begin{equation}
  \begin{gathered}
    f_\mathrm{max} = \mu \cdot f_{z,i} \\
    \mathrm{scale} = \frac{f_\mathrm{max}}{\sqrt{\left(f_{x,i}^\star\right)^2+\left(f_{y,i}^\star\right)^2 + \left(f_\mathrm{max}\right)^2}} \\
  \end{gathered}
\end{equation}

From this, both $f_{x,i}^\star$ and $f_{y,i}^\star$ are scaled by the above quantity, such that the total force generated by the wheel is bounded by the maximum friction cone specified by $\mu$. Finally, the forces are then rotated into the inertial frame and summed across all wheels to get the total force on the car, which is then used in the translational dynamics of the car as shown in equation $(1)$.

The one other aspect of the dynamics that remains is the spin dynamics of each wheel specifically, which we define as a damped rotating disk as follows:

\begin{equation}
  I_w \dot{\omega}_{w,i} = \tau_{\mathrm{cmd},i} - r f_{x,i} - b_w \omega_{w,i}
\end{equation}

Simiarlly, for the passive wheels on the car (i.e. front wheels for RWD), the dynamics are simplified to relax toward the standard kinematic rolling constraint:

\begin{equation}
  \dot{\omega}_{w,i} = \frac{1}{\tau_f}\left(\frac{v_{t,i}}{r}-\omega_{w,i}\right)
\end{equation}

The full state update blends the wheel rotation dynamics via a motor mask, such that $\dot{\omega}_w = \texttt{mask} \cdot \dot{\omega}_\mathrm{drive} + (1- \texttt{mask}) \cdot \dot{\omega}_\mathrm{passive}$.

All combined together, the system of equations that evolve our state \textit{cannot} be formed into a linear state-space model such as $\dot{\mathbf{x}}=A\mathbf{x} + B\mathbf{u}$, due to the nonlinearities in the slip model and the rotation dynamics. However, the system can be linearized around a nominal trajectory, which is what is done in the next section.

\section{Extended Kalman Filter Formulation}

\subsection{Error State Formulation}
While most of the properties of the state $\mathbf{x}$ can be updated via the form $\mathbf{x} \leftarrow \mathbf{x} + K \nu$, we cannot do the same for the orientation $R$, since it evolves on the manifold $SO(3)$. If we ran a standard EKF with a quaternion directly in the state vector, the additive update would push $q$ off of the unit sphere required to properly represent rotations.

%TODO: add reference for error state EKF,
In order to solve this problem, this work formulates a filter in \textbf{error state}. This means that instead of tracking full $SO(3)$, we instead formulate the error orientation as $\delta \boldsymbol{\theta} \in \mathbb{R}^3$, which lives in the tangent space $\mathfrak{so}(3) \cong \mathbb{R}^3$. Doing so allows us to implement a standard state update into the error state EKF, whose state is defined as:  

$$
\mathbf{\delta x} =
\begin{bmatrix}
  \delta \mathbf{p} \in \mathbb{R}^3 \\
  \delta \boldsymbol{\theta} \in \mathbb{R}^3 \\
  \delta \mathbf{v} \in \mathbb{R}^3 \\
  \delta {}^\mathcal{B}\boldsymbol{\omega} \in \mathbb{R}^3 \\
  \delta \boldsymbol{\omega}_w \in \mathbb{R}^4
\end{bmatrix}
$$
where all updates for non-orientation states are standard, except for $\delta \boldsymbol{\theta}$, which is used to update the orientation as follows:

\begin{equation}
  \rotmat{B}{I}+ \leftarrow \rotmat{B}{I} \cdot \exp\left([\delta \boldsymbol{\theta}]_\times\right)
\end{equation}

\subsection{Filter Prediction}

As the components of the system are nonlinear, as stated before we do not have a standard Kalman filter predict step; rather, we propagate the nominal state $\hbf{x}_{k+1|k}$ as:

\begin{equation}
  \hbf{x}_{k+1|k} = f(\hbf{x}_{k|k}, \mathbf{u}_k)
\end{equation}

where $f$ is the integration step of the \mintinline{python}{CarModel} class, implemented in the library. Similarly for the covariance, we do not have a standard $P_{k+1|k} = F P_{k|k} F^T + Q$ update, but rather we linearize the dynamics around the nominal trajectory to get the Jacobian for the error dynamics, as:

\begin{equation}
  F_k = \frac{\partial f}{\partial \delta \mathbf{x}} \biggr|_{\hbf{x}_{k|k}, \mathbf{u}_k}
\end{equation}

And this then feeds into the covariance propagation as:

\begin{equation}
  P_{k+1|k} = F_k P_{k|k} F_k^T + Q
\end{equation}

As $\delta \mathbf{x}$ is reset to zero after each error injection, the prediction error state mean will \textit{always be zero}, as the predict step will only propagate covariance.

\subsection{Filter Update}
